\section{Introduction}\label{sec:intro}

The blended-finance literature has long debated which structures most
effectively mobilize private capital alongside public concessional finance
\citep{convergence2024, attridge2019, lankes2021}. A persistent empirical
limitation runs through that debate: the ``mobilization ratio''---private
dollars moved per public dollar committed---is rarely observed directly at
the project level. Most evidence relies on survey-reported figures or on
coarse program-level aggregates.

The U.S. New Markets Tax Credit offers a setting where the financing
structure behind that ratio becomes visible at the deal level. For every
project, the CDFI Fund's public release records both the subsidized
investment that the intermediary places (the Qualified Low-Income Community
Investment, or QLICI) and the total project cost. I work throughout with
their ratio, which measures how much total project capital accompanies each
dollar of subsidized investment. Section~\ref{sec:data} is careful about
what this ratio is and is not: the government's cost is the credit claimed
against the equity basis, so the ratio describes financing structure rather
than fiscal leverage, and I report both denominators wherever magnitudes
carry weight. The panel covers a single instrument in a single country
across twenty-two years and roughly 8{,}000 projects, which is an unusual
amount of directly measured structure for this literature.

The question this paper asks is simple to state. Projects in
non-metropolitan tracts carry less capital beyond the subsidized investment
than projects in metropolitan ones. Is that gap a \emph{market-structure} phenomenon---rural
projects are fundamentally less leverageable---or an
\emph{intermediary-selection} phenomenon---the Community Development
Entities that deploy in rural tracts are systematically weaker mobilizers
everywhere they operate? The distinction carries direct policy content. A
market-structure story argues for redesigning the credit; a selection story
argues for reallocating it.

In this data the answer is composition, and the layers of the argument
sit in plain view. The raw rural penalty in project-level leverage is
$-0.262$ and precisely estimated. Layering in
origination-year, project-type, and state fixed effects moves it to
$-0.172$. Introducing CDE fixed effects, which compare rural and urban deals
\emph{made by the same intermediary} and draw on the 163 of 343 CDEs that
work both sides of the line, collapses the point estimate to $-0.047$; a
fully nested model retaining state effects alongside CDE effects gives
$-0.060$. This paper is deliberate about what that null does
and does not establish: the mean regression is imprecise (cluster-robust
95\% CI $[-0.245, +0.152]$), and on its own can reject only within-CDE
penalties larger than $0.21$. The precision lives elsewhere, and the paper
reports it where it lives: the median within-CDE penalty is $-0.001$
(SE $0.008$), rural deals within a CDE are no less likely to draw
capital beyond the subsidized investment, and draw no less when they do. What differs
across space is \emph{which} intermediaries show up.

This paper sits inside two literatures. The NMTC evaluation literature
has examined program targeting and neighborhood outcomes extensively
\citep{gurley2009, freedman2012, hargerross2016, freedmankuhns2018,
theodos2020, theodos2022, corinth2024}, with program-evaluation
foundations in \citet{abravanel2013}; the place-based policy literature
supplies the welfare frame \citep{glaesergottlieb2008, klinemoretti2014,
neumarksimpson2015, diamondmcquade2019}. Against both, this paper makes
three contributions, since mobilization has been examined sparingly in
either. First, and subject to the caveat that one cannot prove a negative about
a literature this large, I am not aware of prior econometric work treating
project-level financing structure as the central NMTC outcome, which brings
the mobilization framework into a literature otherwise focused on
neighborhood employment, demographic change, and price effects. Second, it implements a layered
fixed-effects decomposition with CDE fixed effects, separating the
within-CDE rural penalty from the between-CDE selection component---a
decomposition that descriptive work by \citet{theodos2021} motivates but
does not estimate. Third, it brings the excess-mass logic of \citet{chetty2011} and
\citet{kleven2016} to the cross-CDE distribution of non-metro deployment
shares around the program's 20\% target, a test that the target's
prominence in policy discussion would seem to invite and that I have not
found in the published literature. No excess mass appears at the 20\% line,
in deal counts or in dollars, over the full period or over the years after
the proportionality rule entered the code. Compliance plausibly operates
through the allocation agreements that the public release does not
report.

The rest of the paper proceeds as follows. Section~\ref{sec:background}
describes the program's mechanics and the institutional heterogeneity of
the CDE intermediary layer. Section~\ref{sec:data} describes the public
data release and the construction of the leverage outcome.
Section~\ref{sec:strategy} lays out the layered fixed-effects design, its
interpretation, and its limits. Section~\ref{sec:results} presents the
decomposition, heterogeneity, robustness, and the bunching diagnostic.
Section~\ref{sec:discussion} draws out the policy reading, and
Section~\ref{sec:conclusion} concludes. A causal companion design---a
regression-discontinuity analysis at the low-income-community eligibility
threshold---requires a tract-level ACS merge and is deferred to follow-on
work.
